\subsection{Indecomposable Modules and Primordial Modules}

Let A be a ring. Recall the following definition (VII, §4, No.~8, p.~23, Definition 3).

DEFINITION 2. --- An A-module M is called indecomposable if it is not the direct sum of a family of submodules distinct from 0 and M.

By the corollary of Proposition 12 of II, §1, No.~8, p.~209, the following properties are equivalent:

\begin{enumerate}
\def\labelenumi{\alph{enumi})}
\item
  The A-module M is indecomposable.
\item
  The A-module M is nonzero, and every direct factor submodule of M is equal to 0 or M.
\item
  The A-module M is nonzero, and the ring \(\operatorname{End}_{\mathrm{A}}(\mathrm{M})\) contains no idempotent distinct from 0 and \(1_{M}\) .
\end{enumerate}

In particular, since the endomorphism ring of the A-module \(A_{s}\) is isomorphic to the opposite ring of A, we see that the A-module \(A_{s}\) is indecomposable if and only if the ring A is nonzero and its only idempotents are 0 and 1.

Example. --- Suppose that the ring A is a principal ideal domain. The indecomposable finitely generated A-modules are the A-modules isomorphic to either A or \(A/p^{n}A\) , where p is an irreducible element of A and n an integer \textgreater0 (VII, §4, No.~8, p.~24, Proposition 8).

PROPOSITION 3. --- A Noetherian or Artinian A-module M is the direct sum of a finite family of indecomposable submodules.

Let us first prove that every nonzero submodule P of M has an indecomposable direct factor. If this were not the case, then every direct factor submodule of P would be decomposable; proceeding by induction, we could then construct, for every \(n \in N\) , nonzero submodules \(N_{n}^{\prime}\) and \(N_{n}^{\prime\prime}\) of P such that \(P = N_{0}^{\prime} \oplus N_{0}^{\prime\prime}\) and \(N_{n-1}^{\prime} = N_{n}^{\prime} \oplus N_{n}^{\prime\prime}\) for \(n \geqslant 1\) . But then, the sequence of submodules \(N_{0}^{\prime\prime} + \cdots + N_{n}^{\prime\prime}\) would be strictly increasing, and that of the submodules \(N_{n}^{\prime}\) would be strictly decreasing. The module M would be neither Noetherian nor Artinian, contrary to the assumption.

Now, suppose that M is not the direct sum of a finite family of indecomposable submodules. By induction, we construct indecomposable submodules \(P_{n}^{\prime\prime}\) of M and nonzero submodules \(P_{n}^{\prime}\) of M for every \(n \in N\) such that \(M = P_{0}^{\prime} \oplus P_{0}^{\prime\prime}\) and \(P_{n-1}^{\prime} = P_{n}^{\prime} \oplus P_{n}^{\prime\prime}\) for \(n \geqslant 1\) . Indeed, M is nonzero, and the first part of the proof, applied to P = M, provides \(P_{0}^{\prime}\) and \(P_{0}^{\prime\prime}\) . Since the modules \(P_{k}^{\prime}\) and \(P_{k}^{\prime\prime}\) are defined for k \textless{} n, by the first part of the proof, there exist submodules \(P_{n}^{\prime}\) and \(P_{n}^{\prime\prime}\) such that \(P_{n-1}^{\prime} = P_{n}^{\prime} \oplus P_{n}^{\prime\prime}\) with \(P_{n}^{\prime\prime}\) indecomposable. The relation \(M = P_{n}^{\prime} \oplus P_{0}^{\prime\prime} \oplus \cdots \oplus P_{n}^{\prime\prime}\) then implies that \(P_{n}^{\prime} \neq 0\) because M is not the direct sum of a finite family of indecomposable modules.

The sequence of submodules \(P_{0}^{\prime\prime} \oplus \cdots \oplus P_{n}^{\prime\prime}\) is strictly increasing, and the sequence of submodules \(P_{n}^{\prime}\) is strictly decreasing. This contradicts the assumption that M is Noetherian or Artinian.

The question of the uniqueness of the decomposition of a module as a direct sum of indecomposable submodules will be studied in the next subsection.

DEFINITION 3. --- A module is called primordial \(^{(1)}\) if its endomorphism ring is local.

By definition, a local ring is not reduced to 0; consequently, a primordial module is nonzero. Moreover, the A-module \(A_{s}\) is primordial if and only if the ring A is local.

PROPOSITION 4. --- a) A primordial module is indecomposable.

\begin{enumerate}
\def\labelenumi{\alph{enumi})}
\setcounter{enumi}{1}
\tightlist
\item
  An indecomposable module of finite length is primordial.
\end{enumerate}

Let M be an A-module. Suppose that M is primordial; let e be an idempotent in the local ring \(\operatorname{End}_{\mathrm{A}}(\mathrm{M})\) . Since \(e^{2}=e\) , either e is invertible and we have e=1, or 1-e is invertible and we have e=0. This proves that M is indecomposable (VIII, p.~30).

Now, suppose that M is indecomposable and of finite length. By Proposition 2, c) of VIII, p.~27, every endomorphism of M is invertible or nilpotent; the ring \(\operatorname{End}_{\mathrm{A}}(\mathrm{M})\) is therefore local by Example 2 of VIII, p.~26.

The Z-module Z is indecomposable, Noetherian, but not Artinian. Its endomorphism ring is isomorphic to Z, hence is not local. Consequently, Z is not a primordial Z-module.

*Let \(p\) be a prime number. The endomorphism ring of the \(\mathbf{Z}\) -module \(\mathbf{Q}_p / \mathbf{Z}_p\) is isomorphic to the local ring \(\mathbf{Z}_p\) (cf.~VII, §4, p.~65, Exercise 13); it is therefore a primordial \(\mathbf{Z}\) -module.

An injective module is indecomposable if and only if it is primordial (X, §1, n° 9, p.~21, proposition 14).*

